\section{Metodologia de Coleta dos Dados}
Nesta seção, descreve-se o processo utilizado para a coleta dos dados de desempenho dos algoritmos de ordenação. O objetivo foi comparar empiricamente os algoritmos quanto a diversas métricas de eficiência, utilizando vetores de entrada com diferentes tamanhos e ordenações iniciais.

\subsection{Geração dos vetores}
Os vetores utilizados nos testes foram gerados automaticamente por um programa em C, responsável por criar instâncias de entrada com diferentes características. Três tipos de vetores foram considerados:

\begin{enumerate}
    \item Tipo 1 (crescente): vetores já ordenados em ordem crescente.
    \item Tipo 2 (decrescente): vetores ordenados em ordem decrescente.
    \item Tipo 3 (aleatório): vetores com permutação aleatória dos elementos.
\end{enumerate}

Essa variedade permite analisar o comportamento dos algoritmos tanto em casos ideais quanto em situações desfavoráveis.

\subsection{Algoritmos utilizados}
Os algoritmos foram divididos em três categorias:

\begin{enumerate}
    \item Inferiores: algoritmos com complexidade quadrática no pior caso, em geral mais simples e educativos.
    \begin{enumerate}
        \item Bolha, BolhaFlag, Coquetel, Selecao, Insercao.
    \end{enumerate}
    \item Inferiores: algoritmos com complexidade quadrática no pior caso, em geral mais simples e educativos.
    \begin{enumerate}
        \item Shellsort, Mergesort, Quicksort, Quicksort Prob, Quicksort Insercao, Heapsort.
    \end{enumerate}
    \item Lineares: algoritmos com desempenho linear ou quase linear, especialmente eficazes em casos específicos.
    \begin{enumerate}
        \item Contagem, Balde, Radixsort Contagem, Radixsort Balde.
    \end{enumerate}
\end{enumerate}

Cada um dos algoritmos foi implementado e instrumentado para medir métricas internas como número de comparações, movimentos, iterações e tempo de execução.

\subsection{Execuções por amostra}
Para cada combinação de algoritmo, tamanho de vetor e tipo de vetor, foram realizadas dez execuções independentes, com novas instâncias geradas a cada vez. Isso reduz o impacto de variações aleatórias e permite o cálculo de estatísticas como média e desvio padrão, aumentando a confiabilidade da análise.

O número de repetições foi definido com base no equilíbrio entre precisão estatística e tempo de processamento.

\subsection{Execuções por amostra}
Para cada execução, as seguintes métricas foram registradas:

\begin{enumerate}
    \item tempo\_ms: tempo de execução do algoritmo, em milissegundos.
    \item comparacoes: número de comparações realizadas.
    \item movimentos: número de movimentações de elementos (como trocas).
    \item iteracoes: número total de iterações ou chamadas recursivas.
    \item algoritmo, tamanho\_vetor, tipo\_vetor: informações sobre o experimento.
\end{enumerate}

Esses dados foram armazenados em um arquivo CSV para posterior análise.

\subsection{Execuções por amostra}
Todos os testes foram executados em um mesmo ambiente computacional, garantindo consistência nos resultados. O ambiente incluía:

\begin{enumerate}
    \item Sistema operacional: Ubuntu 22.04.
    \item Processador: AMD\textregistered \text{ }Ryzen 5 4600g.
    \item Memória RAM: 16 GB.
    \item Compilador: gcc 11.4.0, com otimização -O2.
    \item Linguagem de programação: C com scripts auxiliares em Bash e Python
\end{enumerate}

A automação foi feita por meio de scripts Bash que executavam o programa com todos os algoritmos, tamanhos e tipos de vetores, salvando os resultados em arquivos CSV.

\section{Gráficos de Desempenho por Tempo de Execução}
A presente seção tem como objetivo comparar o tempo de execução dos algoritmos de ordenação implementados, considerando a classificação adotada neste trabalho: métodos inferiores, métodos superiores e métodos com complexidade linear. Para cada categoria, serão apresentados gráficos com base nas execuções realizadas com vetores de diferentes tamanhos e tipos (aleatório, crescente e decrescente). O eixo x representa o tamanho do vetor, enquanto o eixo y representa o tempo médio de execução, em milissegundos.

A análise visual desses gráficos permite observar o comportamento dos algoritmos conforme o crescimento da entrada, bem como sua sensibilidade à ordenação prévia dos dados.

\subsection{Métodos Inferiores}
Os métodos inferiores incluem os algoritmos de ordenação elementares, cuja complexidade de tempo no pior caso é $O(n^2)$. São eles: Bolha, Bolha com flag, Coquetel, Seleção e Inserção. Estes algoritmos são utilizados principalmente para fins educacionais e em contextos muito específicos, dada sua ineficiência para entradas grandes.

Nos gráficos gerados, observa-se um crescimento acentuado do tempo de execução conforme o tamanho do vetor aumenta, especialmente para os vetores aleatórios e decrescentes. O algoritmo de Inserção apresenta desempenho relativamente melhor em vetores crescentes, como esperado, uma vez que seu número de comparações e movimentos é significativamente reduzido nesse cenário.

[gráfico do tempo por tamanho dos métodos inferiores]

\subsection{Métodos Superiores}
Nesta categoria estão algoritmos com desempenho assintótico superior no caso médio e no pior caso: Shellsort, Mergesort, Quicksort (pivô fixo), Quicksort probabilístico, Quicksort com inserção e Heapsort.

Os gráficos desta seção mostram que o Quicksort com inserção e o Quicksort probabilístico apresentam desempenho bastante eficiente em vetores aleatórios, superando inclusive o Mergesort em algumas situações. O Mergesort, por sua vez, mantém desempenho estável independentemente do tipo de entrada, uma característica esperada de algoritmos com complexidade garantida de O(nlogn).

O Heapsort, apesar de também ser O(nlogn), apresenta desempenho inferior aos demais dessa categoria, o que é coerente com sua maior constante oculta, devido à complexidade das operações de manutenção do heap.

[gráfico do tempo por tamanho dos métodos superiores]

\subsection{Métodos Lineares}
Os métodos com desempenho linear — Contagem, Balde, Radixsort com Contagem e Radixsort com Balde — são indicados para contextos nos quais as chaves são inteiras e com domínio conhecido.

Os gráficos mostram que esses algoritmos superam com folga os demais para vetores de tamanhos maiores, especialmente nos casos aleatórios e crescentes. O Radixsort com Balde apresentou desempenho levemente superior ao Radixsort com Contagem, o que pode ser explicado pela eficiência da sub-rotina utilizada e pela menor complexidade de movimentação dos elementos.

Contagem direta, embora eficiente, apresentou tempo um pouco maior que os algoritmos baseados em Radixsort, o que pode estar relacionado ao custo de inicialização de seus vetores auxiliares.

[gráfico do tempo por tamanho dos métodos lineares]